% Conversão fiel da participação como cliente Hoortech; origem imutável RELATORIO_ATUAL.odt.
\hypertarget{__RefHeading___Toc4470_311817663331}{}

% SOURCE text-h-0034
\refstepcounter{subsection}
\subsection*{6.2.5 Sprint 3}
\label{cliente-hoortech:text-h-0034}
\addcontentsline{toc}{subsection}{6.2.5 Sprint 3}

% SOURCE text-p-0738
A Sprint 3 foi uma grande surpresa e, sem dúvida, um marco no desenvolvimento deste projeto. O time trouxe avanços significativos, e o progresso foi além das expectativas. Primeiramente, o modelo de inteligência artificial estava funcionando perfeitamente e já estava completamente integrado ao front-end. Todas as features planejadas, como a posição da legenda, cores, fontes e até o salvamento do histórico de sinais, foram integradas de forma eficiente. O modelo, que reconhecia as vogais em Libras, estava impecável e completamente funcional.\par

% SOURCE text-p-0739
No início da Sprint, o time enfrentou desafios com o modelo e seu deploy na AWS, principalmente devido ao tamanho do dataset e à quantidade de requisições necessárias para o serviço funcionar corretamente. Após conversas sobre as limitações orçamentárias da AWS e a necessidade de realizar um experimento de POC, definimos que o dataset poderia ser reduzido e que o modelo não precisava ser excessivamente rápido e responsivo. Mesmo com esses limitadores, o time entregou um modelo rápido, com uma requisição a cada 100 milissegundos, e que reconhecia com precisão os sinais presentes no dataset.\par

% SOURCE text-p-0740
O time se comprometeu a expandir o modelo para incluir o alfabeto completo de Libras e, possivelmente, algumas palavras até a Sprint 4. O que me deixou mais contente foi que o modelo estava não apenas funcional, mas também rápido e responsivo. A integração com o front-end foi impecável, e a documentação estava muito bem-feita. Fiquei impressionado ao ver que, com apenas 15 minutos, eu consegui rodar o projeto em minha máquina, seguindo a documentação detalhada que o time preparou. O modelo estava perfeito e funcionando como esperado. Mostrei o resultado para meus professores, pais e colegas, e todos adoraram.\par

% SOURCE text-p-0741
Refletindo sobre o processo, sinto muito orgulho não apenas da equipe, mas também do projeto em si. O que me deixou ainda mais satisfeito foi saber que algo que eu iniciei há mais de um ano, durante uma Hackathona de 24 horas, agora estava sendo levado muito a sério, com grandes avanços sendo alcançados. Ver o time comprometido e entregando resultados tão consistentes foi uma grande satisfação. O projeto estava indo além das minhas expectativas, e, ao ver o modelo funcionando tão bem, senti que o esforço que coloquei no início do projeto havia se transformado em algo real e inovador.\par

